% =========================================================
% CHAPTER 8
% =========================================================

\chapter{Kết luận và hướng phát triển}
\label{Chapter8}

\section{Kết quả đạt được}

Khóa luận đã hoàn thành việc nghiên cứu, thiết kế và xây dựng hệ thống kiểm thử tự động website theo hướng kết hợp giữa mô hình ngôn ngữ lớn, AI agent, browser automation và replay script. So với mục tiêu đã nêu ở Chương~\ref{Chapter1}, hệ thống được thiết kế để hỗ trợ quy trình kiểm thử tự động theo ba trọng tâm chính: giảm công sức tạo test case ban đầu, tăng khả năng tái sử dụng kịch bản kiểm thử và giảm sự phụ thuộc vào LLM trong các lần kiểm thử hồi quy.

\subsection{Về mặt nghiên cứu và kiến thức nền tảng}

Trong quá trình thực hiện đề tài, nhóm đã tìm hiểu các kiến thức nền tảng liên quan đến kiểm thử phần mềm, kiểm thử tự động website, kiểm thử chức năng, kiểm thử hồi quy, browser automation, mô hình ngôn ngữ lớn và AI agent. Các kiến thức này là cơ sở để xác định phạm vi hệ thống, phân tích yêu cầu và lựa chọn hướng thiết kế phù hợp.

Nhóm cũng đã khảo sát các công cụ kiểm thử Web phổ biến như Selenium, Playwright, Cypress và Katalon để làm rõ ưu điểm, hạn chế và vị trí của hệ thống đề xuất. Qua đó, khóa luận xác định hướng tiếp cận chính là không thay thế hoàn toàn các framework kiểm thử truyền thống, mà bổ sung một lớp hỗ trợ giúp người dùng tạo và chạy kiểm thử dễ hơn thông qua mô tả ngôn ngữ tự nhiên.

\subsection{Về mặt xây dựng hệ thống}

Về mặt hiện thực, nhóm đã xây dựng được hệ thống gồm các thành phần chính: frontend, backend API, cơ sở dữ liệu PostgreSQL và agent-worker. Frontend cung cấp giao diện để người dùng quản lý project, test case, dataset, test suite, object repository và report. Backend API xử lý nghiệp vụ, lưu trữ dữ liệu và điều phối yêu cầu chạy test. Agent-worker chịu trách nhiệm chạy test trên trình duyệt thông qua Browser-use và Playwright.

Hệ thống đã hỗ trợ các chức năng chính như:

\begin{itemize}
    \item Đăng ký, đăng nhập và quản lý phiên làm việc của người dùng.
    \item Tạo và quản lý project kiểm thử.
    \item Scan/crawl website để ghi nhận thông tin nền ban đầu.
    \item Tạo test case từ prompt hoặc goal bằng ngôn ngữ tự nhiên.
    \item Lưu trữ test case đã sinh để người dùng có thể chọn lọc, chỉnh sửa và tái sử dụng.
    \item Chạy test case bằng goal-based agent.
    \item Ghi nhận log, screenshot, action history, evidence và verdict.
    \item Sinh replay script từ phiên chạy agent thành công.
    \item Chạy lại replay script bằng Playwright cho các lần kiểm thử hồi quy.
    \item Quản lý test data, dataset và chạy test với nhiều bộ dữ liệu.
    \item Quản lý test suite để chạy nhiều test case theo lô.
    \item Lưu object repository và hỗ trợ locator fallback.
    \item Hỗ trợ phân tích lỗi và hiển thị report sau khi chạy test.
\end{itemize}

Một điểm quan trọng của hệ thống là sự kết hợp giữa hai chế độ thực thi: chạy bằng agent và chạy bằng replay script. Agent được sử dụng chủ yếu trong giai đoạn tạo và khám phá kịch bản ban đầu, khi hệ thống cần suy luận từ mục tiêu kiểm thử bằng ngôn ngữ tự nhiên. Sau khi một phiên chạy thành công, replay script được lưu lại để người dùng có thể chủ động chạy lại bằng Playwright. Cách tổ chức này giúp các kịch bản ổn định có thể được kiểm thử lại mà không cần gọi LLM cho toàn bộ quá trình suy luận hành động, qua đó tạo điều kiện giảm số lần gọi mô hình và tăng tính xác định của luồng thao tác trong kiểm thử hồi quy.

\subsection{Về mặt thực nghiệm và đánh giá}

Về mặt thực nghiệm, hệ thống được đánh giá thông qua case study trên website OrangeHRM. Các kịch bản thực nghiệm tập trung vào những luồng kiểm thử phổ biến của website như đăng nhập, điều hướng, tìm kiếm, chạy test với nhiều bộ dữ liệu và chạy test suite.

Kết quả thực nghiệm cho thấy hệ thống có thể hỗ trợ toàn bộ quy trình kiểm thử từ giai đoạn chuẩn bị project, tạo test case, chạy test, sinh replay script, chạy lại kịch bản và xem báo cáo. Trong 64 batch sinh test case, thời gian sinh trung bình là 3.66 giây mỗi batch và các batch đều tạo được ít nhất một test case có cấu trúc hợp lệ. Kết quả khảo sát cũng cho thấy phần lớn người tham gia chỉ cần chỉnh sửa một phần hoặc rất ít trước khi sử dụng test case được sinh ra. Điều này cho thấy chức năng sinh test case bằng LLM có giá trị như một bản nháp ban đầu, giúp giảm công sức chuẩn bị thủ công cho tester.

Về khả năng tạo cơ sở cho replay script, trong 121 lần Agent initial run trên case study OrangeHRM, hệ thống ghi nhận 110 lần Pass, tương ứng 90.91\%. Các phiên chạy ban đầu thành công này là cơ sở để sinh và lưu replay script cho những kịch bản có thể tái sử dụng trong kiểm thử hồi quy. Khi xét riêng các lần replay thành công, thời gian chạy trung bình là 54.59 giây, thấp hơn agent initial run thành công 14.65\%. Về mức sử dụng LLM, 83 lần chạy agent thuộc phạm vi phân tích token sử dụng tổng cộng 5,349,788 agent tokens, trong khi 46 lần chạy replay script có dữ liệu thực thi và agent token bằng 0. Nếu các lần replay này được chạy lại bằng agent theo mức token trung bình đã ghi nhận, lượng token có thể tránh sử dụng được ước tính khoảng 2.96 triệu tokens. Kết quả này cho thấy replay script có tiềm năng hỗ trợ kiểm thử hồi quy ở các kịch bản đã ổn định, đồng thời giúp giảm số lần cần dùng LLM khi người dùng chọn chạy lại bằng replay script thay vì để agent suy luận lại toàn bộ luồng kiểm thử.

Các bằng chứng như log, screenshot, action history và report giúp người dùng theo dõi quá trình thực thi và phân tích nguyên nhân khi test thất bại. Đây là phần hỗ trợ quan trọng để tester quyết định kịch bản nào có thể tái sử dụng trực tiếp, kịch bản nào cần chỉnh sửa test data, locator, assertion hoặc chạy lại bằng agent.

Tuy nhiên, kết quả thực nghiệm cũng cho thấy hiệu quả của hệ thống phụ thuộc vào chất lượng prompt, trạng thái website mục tiêu, dữ liệu kiểm thử, khả năng tìm lại phần tử giao diện và độ ổn định của mô hình AI. Đây là những yếu tố cần tiếp tục được cải thiện trong các phiên bản sau.

\subsection{Đối chiếu với mục tiêu đề tài}

Bảng~\ref{tab:chapter8_objective_alignment} tổng hợp mức độ đáp ứng của hệ thống đối với các mục tiêu chính đã nêu ở Chương~\ref{Chapter1}.

\begin{table}[H]
\centering
\caption{Đối chiếu kết quả đạt được với mục tiêu đề tài}
\label{tab:chapter8_objective_alignment}
\small
\renewcommand{\arraystretch}{1.2}
\begin{tabular}{|L{4.4cm}|L{5.2cm}|L{4.6cm}|}
\hline
\textbf{Mục tiêu} & \textbf{Cách hệ thống đáp ứng} & \textbf{Kết quả ghi nhận} \\
\hline
Giảm công sức tạo test case, test script và test data & Dùng LLM để sinh test case/dataset từ mô tả tự nhiên; dùng agent để tạo luồng chạy ban đầu và sinh replay script từ phiên chạy thành công. & 64 batch sinh test case có cấu trúc hợp lệ; thời gian sinh test case trung bình 3.66 giây; sinh dataset trung bình 2.57 giây. \\
\hline
Tái sử dụng test case và kịch bản kiểm thử & Lưu test case, test data, test suite, replay script, object repository và evidence để dùng lại trong các lần kiểm thử sau. & Agent initial run đạt 110/121 lần Pass, tương ứng 90.91\%; đây là cơ sở để sinh replay script cho các kịch bản có thể tái sử dụng trong kiểm thử hồi quy. \\
\hline
Giảm chi phí sử dụng LLM trong kiểm thử hồi quy & Sau lần chạy thành công bằng agent, replay script được lưu để người dùng có thể chủ động chạy lại các kịch bản ổn định bằng Playwright thay vì chạy lại bằng agent. & 83 lần chạy agent thuộc phạm vi phân tích token sử dụng 5,349,788 tokens; 46 replay run có dữ liệu thực thi và agent token bằng 0, tương ứng khoảng 2.96 triệu tokens có thể tránh sử dụng theo ước tính. \\
\hline
Hỗ trợ phân tích và bảo trì kiểm thử & Ghi nhận log, screenshot, action history, verdict, failure type và report để truy vết lỗi. & Report/evidence giúp tester xác định lỗi do locator, dữ liệu, assertion, môi trường hoặc prompt chưa rõ. \\
\hline
\end{tabular}
\end{table}

Như vậy, kết quả của đề tài phù hợp với định hướng ban đầu: LLM và agent được dùng để giảm công sức ở giai đoạn tạo kịch bản đầu tiên, còn replay script, test suite, dataset và object repository được dùng để tăng khả năng tái sử dụng trong các lần kiểm thử tiếp theo. Hệ thống đã ghi nhận số token sử dụng trong quá trình sinh test case, sinh dataset và thực thi bằng agent; từ số token này, chi phí sử dụng LLM có thể được ước tính theo đơn giá của mô hình được cấu hình.

\section{Hạn chế của hệ thống}

Mặc dù hệ thống đã hoàn thành các chức năng chính và đáp ứng được các mục tiêu ban đầu của đề tài, vẫn còn một số giới hạn cần được ghi nhận để định hướng cho các phiên bản phát triển tiếp theo.

Thứ nhất, kết quả thực thi bằng agent vẫn phụ thuộc vào chất lượng prompt và trạng thái giao diện tại thời điểm chạy. Khi prompt mô tả chưa đủ rõ mục tiêu kiểm thử, dữ liệu đầu vào hoặc điều kiện kỳ vọng, agent có thể hiểu sai nhiệm vụ hoặc chọn hành động chưa phù hợp. Ngoài ra, với các website có nhiều trạng thái giao diện động, popup, dữ liệu tải chậm hoặc luồng điều hướng phức tạp, quá trình suy luận và thao tác của agent có thể chưa ổn định trong mọi trường hợp.

Thứ hai, khả năng tái sử dụng replay script vẫn phụ thuộc vào mức độ ổn định của giao diện website mục tiêu. Replay script có hiệu quả tốt hơn đối với các luồng nghiệp vụ đã ổn định, trong đó locator, dữ liệu đầu vào và điều kiện kiểm tra không thay đổi quá nhiều. Khi website thay đổi mạnh về cấu trúc DOM, thay đổi luồng nghiệp vụ, xóa phần tử hoặc thay đổi logic xử lý, người dùng vẫn cần xem lại report và cập nhật test case hoặc replay script. Vì vậy, khả năng tái sử dụng trong hệ thống được hiểu là tái sử dụng có điều kiện, phù hợp với các luồng nghiệp vụ còn ổn định qua các phiên bản website.

Thứ ba, thực nghiệm hiện tại chủ yếu đánh giá khả năng tạo và chạy lại replay script trên cùng website mục tiêu. Hệ thống chưa được đánh giá theo một quy trình kiểm soát gồm nhiều phiên bản liên tiếp của cùng website, chẳng hạn chạy trên phiên bản thứ nhất, thay đổi sang phiên bản thứ hai, chạy lại cùng replay script và đo tỷ lệ script sử dụng được trực tiếp, cần chỉnh sửa hoặc phải tạo lại. Vì vậy, khả năng tái sử dụng qua thay đổi phiên bản mới chỉ được phản ánh gián tiếp, chưa được đo lường đầy đủ.

Thứ tư, cơ chế ghi nhận object repository còn phụ thuộc vào thông tin định danh của phần tử tại thời điểm quan sát. Trong một số trường hợp, cùng một phần tử giao diện như nút bấm hoặc ô nhập liệu có thể không thay đổi về mặt hiển thị, nhưng cấu trúc DOM, XPath, selector, accessible name hoặc các thuộc tính nhận dạng bên dưới đã thay đổi giữa các lần chạy. Khi đó, hệ thống có thể ghi nhận phần tử này như một object mới thay vì ánh xạ về object đã có. Điều này có thể làm tăng số lượng object trùng nghĩa trong repository và gây khó khăn hơn cho việc bảo trì locator về lâu dài.

Thứ năm, trong phạm vi khóa luận, hệ thống được thiết kế và đánh giá tập trung vào kiểm thử chức năng trên giao diện Web. Các loại kiểm thử khác như kiểm thử hiệu năng, kiểm thử bảo mật, kiểm thử API độc lập hoặc kiểm thử ứng dụng di động chưa được triển khai và đánh giá sâu. Đây không phải là mục tiêu chính của đề tài hiện tại, nhưng là các hướng có thể mở rộng trong tương lai nếu hệ thống được phát triển thành một nền tảng kiểm thử toàn diện hơn.

Thứ sáu, chức năng phân tích lỗi bằng AI hiện đóng vai trò hỗ trợ người dùng giải thích nguyên nhân thất bại và đưa ra gợi ý ban đầu. Hệ thống chưa tự động sửa mọi lỗi hoặc tự cập nhật test case, assertion, dữ liệu kiểm thử hay replay script mà không có sự xác nhận của người dùng. Cách tiếp cận này giúp đảm bảo người dùng vẫn giữ vai trò kiểm soát cuối cùng trong quá trình bảo trì kịch bản kiểm thử.

\section{Bài học kinh nghiệm}

Trong quá trình thực hiện đề tài, nhóm rút ra một số bài học kinh nghiệm quan trọng.

Trước hết, đối với một hệ thống có nhiều thành phần như frontend, backend, database, worker, browser automation và LLM, việc thiết kế kiến trúc tách biệt ngay từ đầu là rất cần thiết. Cách tách thành phần giúp quá trình phát triển, kiểm thử và sửa lỗi dễ kiểm soát hơn.

Thứ hai, các chức năng liên quan đến kiểm thử tự động trên trình duyệt cần có cơ chế ghi nhận log và evidence đầy đủ. Khi test thất bại, nếu không có screenshot, action history, URL hiện tại hoặc error message, việc xác định nguyên nhân lỗi sẽ rất khó khăn.

Thứ ba, không nên phụ thuộc hoàn toàn vào AI agent trong mọi lần chạy. Agent phù hợp với giai đoạn tạo và khám phá luồng kiểm thử ban đầu, trong khi replay script có thể hỗ trợ kiểm thử hồi quy với số lần gọi LLM ít hơn và tính xác định cao hơn về luồng thao tác đối với các kịch bản đã ổn định.

Thứ tư, dữ liệu kiểm thử, trạng thái website và locator có ảnh hưởng lớn đến kết quả chạy test. Vì vậy, hệ thống cần có cơ chế quản lý dataset, object repository và fallback locator để hỗ trợ bảo trì kịch bản kiểm thử.

Cuối cùng, việc viết báo cáo cần đi song song với quá trình phát triển hệ thống. Khi các quyết định thiết kế, khó khăn kỹ thuật và kết quả thực nghiệm được ghi nhận kịp thời, nội dung báo cáo sẽ đầy đủ và phản ánh đúng quá trình thực hiện đề tài hơn.

\section{Hướng phát triển trong tương lai}

Trong tương lai, hệ thống có thể được mở rộng theo nhiều hướng khác nhau.

Thứ nhất, hệ thống có thể cải thiện khả năng sinh test case bằng cách bổ sung nhiều ngữ cảnh hơn cho mô hình AI, chẳng hạn như kết quả scan/crawl, mô tả nghiệp vụ, trạng thái DOM, screenshot và lịch sử các lần chạy trước. Điều này có thể giúp test case được sinh ra rõ ràng và phù hợp hơn với website mục tiêu, từ đó giảm thêm công sức chỉnh sửa của tester sau khi sinh.

Thứ hai, cơ chế tái sử dụng cần được phát triển sâu hơn ở cả mức test case và replay script. Hệ thống có thể bổ sung chức năng versioning cho test case, so sánh thay đổi giữa các phiên bản website, gợi ý test case bị ảnh hưởng và cho phép tester cập nhật một phần thay vì tạo lại toàn bộ kịch bản.

Thứ ba, cơ chế replay script có thể được mở rộng với nhiều loại action và assertion hơn. Ngoài các thao tác cơ bản như navigate, fill và click, hệ thống có thể hỗ trợ thêm các thao tác phức tạp như upload file, drag-and-drop, chọn ngày, xử lý modal, kiểm tra bảng dữ liệu hoặc kiểm tra trạng thái API phía sau.

Thứ tư, object repository và self-healing có thể được cải tiến để tăng khả năng tìm lại phần tử khi giao diện thay đổi. Hệ thống có thể bổ sung cơ chế chấm điểm locator, học từ các lần fallback thành công, nhận diện các object trùng nghĩa dù thông tin định danh thay đổi hoặc đề xuất cập nhật object sau khi replay script gặp lỗi.

Thứ năm, chức năng AI analysis có thể được phát triển sâu hơn. Thay vì chỉ gợi ý nguyên nhân lỗi, hệ thống có thể đề xuất bản sửa cụ thể cho step, assertion, selector hoặc dữ liệu kiểm thử, sau đó cho phép người dùng xem lại và xác nhận trước khi lưu.

Thứ sáu, hệ thống có thể mở rộng khả năng thực nghiệm và đánh giá trên nhiều website hơn. Việc kiểm thử trên nhiều loại website khác nhau như thương mại điện tử, quản trị nội bộ, form nghiệp vụ hoặc dashboard dữ liệu sẽ giúp đánh giá rõ hơn độ ổn định và khả năng áp dụng thực tế của hệ thống.

Thứ bảy, hệ thống có thể được triển khai theo hướng hỗ trợ nhiều worker hoặc hàng đợi thực thi để chạy nhiều test song song. Điều này giúp hệ thống phù hợp hơn với nhu cầu kiểm thử hồi quy lớn hoặc chạy test suite định kỳ.

Cuối cùng, hệ thống có thể bổ sung các cơ chế bảo mật nâng cao hơn như mã hóa thông tin xác thực website, che dữ liệu nhạy cảm trong screenshot/log, phân quyền theo nhóm người dùng và tích hợp với các hệ thống CI/CD. Đây là những hướng cần thiết nếu hệ thống được phát triển thành một nền tảng kiểm thử tự động có khả năng ứng dụng trong môi trường thực tế.

\section{Kết luận chung}

Khóa luận đã đề xuất và xây dựng một hệ thống kiểm thử tự động website theo hướng kết hợp giữa LLM, AI agent, browser automation và replay script. Hệ thống cho phép người dùng mô tả mục tiêu kiểm thử bằng ngôn ngữ tự nhiên, sinh test case/test data, chạy test trên trình duyệt, ghi nhận bằng chứng, sinh replay script và hỗ trợ chạy lại kịch bản kiểm thử.

Kết quả đạt được cho thấy hướng tiếp cận này phù hợp với mục tiêu ban đầu của đề tài. LLM giúp tester tạo bản nháp test case và dataset nhanh hơn so với chuẩn bị thủ công hoàn toàn. Agent hỗ trợ khám phá và hiện thực phiên chạy đầu tiên trên trình duyệt thật. Replay script, test suite, dataset và object repository tạo nền tảng để tái sử dụng các kịch bản đã ổn định trong kiểm thử hồi quy, từ đó giảm nhu cầu gọi LLM lặp lại ở các lần chạy sau. Mặc dù hệ thống vẫn còn một số hạn chế về độ ổn định của replay script, khả năng phục hồi khi giao diện thay đổi và phạm vi thực nghiệm, đề tài đã xây dựng được một nền tảng ban đầu cho quy trình kiểm thử Web có hỗ trợ AI nhưng không phụ thuộc hoàn toàn vào AI trong toàn bộ vòng đời kiểm thử.
